\documentclass[a4paper,12pt]{article}

\usepackage[T2A]{fontenc}
\usepackage[utf8]{inputenc}
\usepackage[russian]{babel}

\usepackage{amsmath,amssymb,amsthm,mathtools}
\usepackage{helvet}
\renewcommand{\familydefault}{\sfdefault}

\usepackage{icomma}
\usepackage[a4paper,left=20mm,right=20mm,top=20mm,bottom=22mm]{geometry}
\usepackage{microtype}

\usepackage{xcolor}
\definecolor{accent}{HTML}{008080}
\definecolor{darkaccent}{HTML}{004D4D}
\definecolor{textgray}{HTML}{333333}
\definecolor{softgray}{HTML}{6A7373}
\definecolor{lightaccent}{HTML}{DCEEEE}

\usepackage{tikz}
\usetikzlibrary{calc}

\usepackage{tkz-euclide}

\usepackage{pgfplots}
\pgfplotsset{compat=1.18}

\usepackage{tabularx,booktabs,longtable}
\usepackage{enumitem}
\usepackage{parskip}
\usepackage{fancyhdr}
\usepackage{setspace}

\usepackage[
    unicode,
    colorlinks=true,
    linkcolor=darkaccent,
    urlcolor=darkaccent
]{hyperref}

\onehalfspacing
\setlength{\parindent}{0pt}
\setlength{\headheight}{26pt}

\setlist[itemize]{
    leftmargin=1.6em,
    itemsep=0.3em,
    topsep=0.35em
}

\setlist[enumerate]{
    leftmargin=1.9em,
    itemsep=0.55em,
    topsep=0.45em
}

\pagestyle{fancy}
\fancyhf{}
\fancyhead[L]{\small\color{softgray}ОГЭ по математике}
\fancyhead[R]{\small\color{softgray}Ярослав Верещагин}
\fancyfoot[C]{\small\color{softgray}\thepage}

\renewcommand{\headrulewidth}{0.4pt}
\renewcommand{\headrule}{%
    \hbox to\headwidth{%
        \color{lightaccent}%
        \leaders\hrule height \headrulewidth\hfill
    }%
}

\newcommand{\topic}[1]{%
    \par\vspace{0.9em}%
    {\Large\bfseries\color{darkaccent}#1}\par
    \vspace{0.25em}
    \noindent
    \color{accent}\rule{\linewidth}{0.6pt}
    \color{black}
    \vspace{0.55em}
}

\newcommand{\subtopic}[1]{%
    \par\vspace{0.55em}%
    {\large\bfseries\color{textgray}#1}\par
    \vspace{0.2em}
}

\newcommand{\important}[1]{%
    \textbf{\color{darkaccent}#1}%
}

\begin{document}

% =========================================================
% ТИТУЛЬНЫЙ ЛИСТ
% =========================================================

\begin{titlepage}
\thispagestyle{empty}
\centering

\vspace*{22mm}

{\Large\color{softgray}ОГЭ по математике\par}

\vspace{12mm}

{\fontsize{32}{38}\selectfont
\bfseries
\color{darkaccent}
Чек-лист
\par}

\vspace{5mm}

{\LARGE\bfseries
Арифметика дробей, числовая прямая
\par}

\vspace{2mm}

{\LARGE\bfseries
и оценка чисел
\par}

\vspace{18mm}

{\large
Повторение базовых вычислительных приёмов
\par}

{\large
в формате экзаменационных задач
\par}

\vfill

{\large\bfseries
Лёвин Артём Александрович
\par}

\vspace{2mm}

{\large
эксклюзивно для Ярослава Верещагина
\par}

\vspace{8mm}

{\large
\today
\par}

\vspace{16mm}

\begin{tikzpicture}[x=1cm,y=1cm]
    \draw[
        accent!55,
        line width=1.2pt
    ]
    (-7.6,0)
    .. controls (-4.7,1.05) and (-2.2,-0.65) ..
    (0,0.15)
    .. controls (2.3,0.95) and (4.6,-0.55) ..
    (7.6,0.15);
\end{tikzpicture}

\end{titlepage}

\setcounter{page}{1}

% =========================================================
% 1
% =========================================================

\topic{1. Обыкновенные дроби: сложение и вычитание}

Для двух дробей можно сразу использовать перекрёстное умножение
(часто его называют \emph{«методом улыбочки»}):

\[
\frac{a}{b}\pm\frac{c}{d}
=
\frac{ad\pm bc}{bd},
\qquad
b\ne0,\quad d\ne0.
\]

Этот приём особенно удобен, когда знаменатели небольшие и нет
необходимости отдельно искать наименьший общий знаменатель.

\subtopic{Пример}

\[
\frac{4}{25}+\frac{15}{4}
=
\frac{4\cdot4+15\cdot25}{25\cdot4}
=
\frac{16+375}{100}
=
\frac{391}{100}
=
3{,}91.
\]

\important{Проверка.}
Перед окончательным ответом посмотрите, можно ли сократить дробь.

Если знаменатель равен $10$, $100$, $1000$ и т.~п.,
удобно сразу перейти к десятичной записи.

\subtopic{Что важно не перепутать}

При вычитании знак между перекрёстными произведениями также остаётся
знаком вычитания:

\[
\frac{a}{b}-\frac{c}{d}
=
\frac{ad-bc}{bd}.
\]

Особенно внимательно следите за знаками, если одна из дробей
отрицательная.

% =========================================================
% 2
% =========================================================

\topic{2. Деление обыкновенных дробей}

Деление заменяем умножением на дробь, обратную
\emph{делителю}:

\[
\frac{a}{b}:\frac{c}{d}
=
\frac{a}{b}\cdot\frac{d}{c},
\]

причём

\[
b\ne0,\qquad d\ne0,\qquad c\ne0.
\]

Последнее условие означает, что делитель
$\dfrac{c}{d}$ не должен быть равен нулю.

\subtopic{Алгоритм}

\begin{enumerate}[label=\arabic*)]
    \item замените знак деления на умножение;
    \item переверните только дробь, стоящую после знака деления;
    \item выполните сокращение крест-накрест, если оно возможно;
    \item перемножьте оставшиеся числители и знаменатели;
    \item сократите результат, если это возможно.
\end{enumerate}

\subtopic{Пример}

\[
\frac{12}{5}:\frac{15}{2}
=
\frac{12}{5}\cdot\frac{2}{15}.
\]

До умножения удобно выполнить сокращение:

\[
\frac{12}{15}=\frac45.
\]

Поэтому

\[
\frac{12}{5}\cdot\frac{2}{15}
=
\frac{4\cdot2}{5\cdot5}
=
\frac{8}{25}.
\]

Знаменатель $25$ удобно превратить в $100$:

\[
\frac{8}{25}
=
\frac{8\cdot4}{25\cdot4}
=
\frac{32}{100}
=
0{,}32.
\]

\important{Типичная ошибка.}
Перевернуть первую дробь или обе дроби.

Переворачивается только делитель.

% =========================================================
% 3
% =========================================================

\topic{3. Смешанные числа}

Перед вычислениями смешанные числа удобно переводить
в неправильные дроби.

Общее правило:

\[
a\frac{b}{c}
=
\frac{ac+b}{c},
\qquad
c\ne0.
\]

\subtopic{Пример 1}

\[
2\frac35
=
\frac{2\cdot5+3}{5}
=
\frac{13}{5}.
\]

При необходимости:

\[
\frac{13}{5}=2{,}6.
\]

\subtopic{Пример 2}

\[
5\frac13
=
\frac{5\cdot3+1}{3}
=
\frac{16}{3}.
\]

\important{Практический приём.}
Перед умножением, делением или длинной цепочкой вычислений
смешанное число почти всегда полезно сначала превратить
в неправильную дробь.

% =========================================================
% 4
% =========================================================

\topic{4. Десятичные дроби: сложение и вычитание}

При сложении и вычитании столбиком действует главное правило:

\begin{center}
\important{запятая должна находиться под запятой.}
\end{center}

Целые записываем под целыми, десятые под десятыми,
сотые под сотыми.

Недостающие разряды справа можно дополнить нулями.

Например,

\[
2{,}3400+3{,}15
=
2{,}3400+3{,}1500
=
5{,}4900
=
5{,}49.
\]

Дописывание нулей справа не меняет значение числа:

\[
3{,}15=3{,}150=3{,}1500.
\]

\important{Мнемоника.}
При сложении и вычитании --- \textbf{запятая под запятой}.

% =========================================================
% 5
% =========================================================

\topic{5. Умножение десятичных дробей}

При умножении столбиком числа выравнивают
\important{по правому краю}.

Сначала временно работаем с числами как с целыми.

После умножения возвращаем десятичную запятую.

\subtopic{Алгоритм}

\begin{enumerate}[label=\arabic*)]
    \item посчитайте общее количество цифр после запятой
    в обоих множителях;
    \item временно не учитывайте запятые;
    \item выполните обычное умножение;
    \item в результате отделите справа столько цифр,
    сколько получилось на первом шаге.
\end{enumerate}

\subtopic{Пример}

\[
2{,}34\cdot2{,}1.
\]

В первом множителе две цифры после запятой,
во втором --- одна.

Всего три цифры.

Как целые числа:

\[
234\cdot21=4914.
\]

Возвращаем три знака после запятой:

\[
2{,}34\cdot2{,}1=4{,}914.
\]

\important{Не путайте два правила.}

\begin{itemize}
    \item сложение и вычитание:
    \textbf{запятая под запятой};
    \item умножение:
    \textbf{выравнивание по правому краю}.
\end{itemize}

% =========================================================
% 6
% =========================================================

\topic{6. Деление десятичных дробей}

При делении важно убрать запятую в
\important{делителе}.

Для этого делимое и делитель умножают
на одну и ту же степень числа $10$.

\subtopic{Пример 1}

\[
1{,}2:0{,}5.
\]

Чтобы число $0{,}5$ превратить в целое, умножаем
оба числа на $10$:

\[
1{,}2:0{,}5
=
12:5
=
2{,}4.
\]

\subtopic{Пример 2}

\[
4{,}2:0{,}07.
\]

У делителя две цифры после запятой,
поэтому умножаем оба числа на $100$:

\[
4{,}2:0{,}07
=
420:7
=
60.
\]

\subtopic{Пример 3}

\[
0{,}05:1{,}2.
\]

Для удаления запятой именно в делителе достаточно умножить
числа на $10$:

\[
0{,}05:1{,}2
=
0{,}5:12.
\]

Можно пойти дальше и убрать запятые вообще:

\[
0{,}05:1{,}2
=
5:120
=
\frac{1}{24}.
\]

\important{Главное ограничение.}
Делить на ноль нельзя.

Если делитель равен $0$, выражение не определено.

% =========================================================
% 7
% =========================================================

\topic{7. Полезные дроби, которые стоит узнавать сразу}

Некоторые соответствия удобно помнить без вычислений.

\begin{table}[h]
\centering
\caption{Ходовые соответствия обычных и десятичных дробей}
\renewcommand{\arraystretch}{1.35}

\begin{tabular}{@{}cccccc@{}}
\toprule
Дробь
&
$\dfrac12$
&
$\dfrac14$
&
$\dfrac15$
&
$\dfrac18$
&
$\dfrac13$
\\
\midrule
Десятичная запись
&
$0{,}5$
&
$0{,}25$
&
$0{,}2$
&
$0{,}125$
&
$0{,}(3)$
\\
\bottomrule
\end{tabular}
\end{table}

Также полезно помнить:

\[
\frac16=0{,}1(6),
\qquad
\frac23=0{,}(6).
\]

Периодическую дробь можно получить обычным делением столбиком.

Для сравнения чисел обычно достаточно нескольких первых
знаков после запятой.

% =========================================================
% 8
% =========================================================

\topic{8. Сравнение дробей}

Для сравнения обыкновенных дробей особенно полезны два метода.

\subtopic{Способ 1. Перевод в десятичные дроби}

Если варианты ответа уже даны десятичными дробями,
этот способ обычно самый быстрый.

Например,

\[
\frac16\approx0{,}1667,
\qquad
\frac14=0{,}25.
\]

Следовательно,

\[
\frac16<0{,}2<\frac14.
\]

\subtopic{Пример с несколькими дробями}

Пусть требуется определить положение числа $\dfrac38$ среди дробей

\[
\frac16,\qquad
\frac13,\qquad
\frac12,\qquad
\frac23.
\]

Переведём нужные числа в десятичную форму:

\[
\frac38=0{,}375,
\]

\[
\frac13\approx0{,}333,
\qquad
\frac12=0{,}5.
\]

Поэтому

\[
\frac13<\frac38<\frac12.
\]

\subtopic{Способ 2. Общий знаменатель}

Можно привести все дроби к одинаковому знаменателю.

Например,

\[
\frac38=\frac9{24},
\qquad
\frac16=\frac4{24},
\qquad
\frac13=\frac8{24},
\qquad
\frac12=\frac{12}{24}.
\]

После этого достаточно сравнить числители:

\[
8<9<12.
\]

Значит,

\[
\frac13<\frac38<\frac12.
\]

\subtopic{Как выбрать способ}

\begin{itemize}
    \item если общий знаменатель виден сразу,
    можно использовать его;
    \item если знаменатели $6$, $7$, $9$, $12$ и т.~п.,
    десятичные приближения часто экономят время;
    \item если ответы уже записаны десятичными дробями,
    обычно выгодно перевести все границы именно в этот вид.
\end{itemize}

% =========================================================
% 9
% =========================================================

\topic{9. Числовая прямая и приближённая оценка}

Если число задано положением точки на числовой прямой,
точное значение часто искать не требуется.

Можно выбрать удобное приближённое число.

Важно сохранить:

\begin{itemize}
    \item знак;
    \item порядок точек слева направо;
    \item положение относительно нуля;
    \item сравнительные расстояния от нуля;
    \item соотношения модулей, если они видны на рисунке.
\end{itemize}

Например, для схемы

\begin{center}
\begin{tikzpicture}[x=1.7cm,y=1cm]

    \draw[->,line width=0.9pt]
    (-2.8,0)--(2.8,0)
    node[right] {$x$};

    \foreach \x in {-2,-1,0,1,2}{
        \draw
        (\x,0.09)--(\x,-0.09);
    }

    \fill[accent]
    (-2,0)
    circle (2.2pt)
    node[above=5pt] {$A$};

    \fill[accent]
    (-1,0)
    circle (2.2pt)
    node[above=5pt] {$B$};

    \fill[accent]
    (1,0)
    circle (2.2pt)
    node[above=5pt] {$C$};

    \node[below=5pt] at (0,0) {$0$};

\end{tikzpicture}
\end{center}

удобно взять

\[
A=-2,\qquad B=-1,\qquad C=1.
\]

При этом

\[
|B|=|C|=1.
\]

Это связано с тем, что модуль числа показывает расстояние
от соответствующей точки до нуля.

% =========================================================
% 10
% =========================================================

\topic{10. Подстановка отрицательных чисел}

Если вместо переменной подставляется отрицательное число,
его безопаснее брать в скобки.

Например,

\[
2-a
\]

при

\[
a=-4{,}4
\]

даёт

\[
2-(-4{,}4)
=
2+4{,}4
=
6{,}4.
\]

Скобки особенно важны при умножении и возведении в степень:

\[
(-3)^2=9,
\]

но

\[
-3^2=-9.
\]

Во втором выражении сначала вычисляется степень:

\[
3^2=9,
\]

после чего применяется минус.

\important{Рабочее правило.}
При подстановке отрицательного числа сначала поставьте скобки,
а уже потом выполняйте действия.

% =========================================================
% 11
% =========================================================

\topic{11. Минус в дроби}

Для ненулевого $b$ справедливо:

\[
\frac{a}{-b}
=
-\frac{a}{b}
=
\frac{-a}{b}.
\]

Минус можно:

\begin{itemize}
    \item оставить в знаменателе;
    \item перенести в числитель;
    \item вынести перед всей дробью.
\end{itemize}

Например,

\[
\frac{10}{-28}
=
-\frac{10}{28}
=
-\frac{5}{14}.
\]

Это отрицательное число.

\important{Типичная ошибка.}
Минус нельзя просто убрать.

\subtopic{Отдельно о неравенствах}

Если обе части неравенства умножаются или делятся
на отрицательное число, знак неравенства меняется:

\[
-2x>6
\]

\[
x<-3.
\]

Это правило относится именно к преобразованию неравенства.

% =========================================================
% 12
% =========================================================

\topic{12. Проверка утверждений о числе на числовой прямой}

Пусть по рисунку видно, что

\[
a\approx-4{,}4.
\]

Чтобы проверить утверждение, можно подставить это приближённое
значение в выражение.

Например:

\[
a+4\approx-4{,}4+4=-0{,}4.
\]

Следовательно,

\[
a+4>0
\]

неверно.

Далее:

\[
a+5\approx-4{,}4+5=0{,}6.
\]

Поэтому

\[
a+5<0
\]

также неверно.

Для выражения

\[
2-a
\]

получаем

\[
2-(-4{,}4)=6{,}4>0.
\]

Значит, утверждение

\[
2-a>0
\]

верно.

\important{Идея метода.}
Здесь не нужна высокая точность.

Достаточно выбрать число, которое правильно отражает
положение точки на прямой.

% =========================================================
% 13
% =========================================================

\topic{13. Оценка квадратных корней}

Чтобы понять, между какими соседними целыми числами находится
$\sqrt n$, найдите ближайшие полные квадраты.

Если

\[
k^2<n<(k+1)^2,
\]

то

\[
k<\sqrt n<k+1.
\]

Здесь предполагается

\[
n\ge0.
\]

\subtopic{Пример}

Для $\sqrt{53}$:

\[
49<53<64.
\]

Но

\[
49=7^2,
\qquad
64=8^2.
\]

Следовательно,

\[
7<\sqrt{53}<8.
\]

\begin{center}
\begin{tikzpicture}[x=1.6cm,y=1cm]

    \draw[->,line width=0.9pt]
    (6.5,0)--(8.5,0);

    \draw
    (7,0.09)--(7,-0.09)
    node[below=4pt] {$7$};

    \draw
    (8,0.09)--(8,-0.09)
    node[below=4pt] {$8$};

    \fill[accent]
    (7.28,0)
    circle (2.2pt)
    node[above=5pt] {$\sqrt{53}$};

\end{tikzpicture}
\end{center}

Аналогично,

\[
36<42<49
\]

даёт

\[
6<\sqrt{42}<7.
\]

\important{Экзаменационный приём.}
Если требуется только указать интервал,
нет необходимости вычислять корень с высокой точностью.

% =========================================================
% 14
% =========================================================

\topic{14. Экзаменационная стратегия сравнения чисел}

При задаче вида

\begin{center}
«Какое число расположено между двумя заданными числами?»
\end{center}

полезен следующий алгоритм.

\begin{enumerate}[label=\arabic*)]
    \item Определите, в каком виде записаны варианты ответа.
    \item Приведите границы к тому же удобному виду.
    \item Если используются дроби, при необходимости получите
    десятичные приближения.
    \item Сравнивайте числа последовательно.
    \item Проверьте обе границы:
    \[
    \text{левая граница}
    <
    \text{ответ}
    <
    \text{правая граница}.
    \]
\end{enumerate}

Например, если нужно проверить число $\dfrac59$ относительно
границ $\dfrac37$ и $\dfrac47$, оценим:

\[
\frac37\approx0{,}4286,
\]

\[
\frac59\approx0{,}5556,
\]

\[
\frac47\approx0{,}5714.
\]

Поэтому

\[
\frac37<\frac59<\frac47.
\]

% =========================================================
% 15
% =========================================================

\topic{15. Итоговый чек-лист самопроверки}

Перед тем как отметить ответ в экзаменационной работе,
быстро проверьте следующее.

\begin{enumerate}[label=\arabic*)]

    \item Верно ли переписан знак действия?

    \item При сложении или вычитании дробей правильно ли
    выполнены перекрёстные произведения?

    \item При делении дробей перевёрнут ли именно делитель?

    \item Не выполняется ли деление на ноль?

    \item Сокращены ли дроби там, где это возможно?

    \item При сложении десятичных дробей стоит ли
    запятая под запятой?

    \item При умножении десятичных дробей правильно ли
    посчитано количество знаков после запятой?

    \item При делении десятичных дробей убрана ли запятая
    именно в делителе?

    \item При подстановке отрицательного числа поставлены ли
    скобки?

    \item Не потерян ли минус у дроби?

    \item Если преобразовывалось неравенство умножением или
    делением на отрицательное число, изменён ли его знак?

    \item При сравнении дробей приведены ли числа
    к удобному общему виду?

    \item Для квадратного корня найдены ли соседние
    полные квадраты?

\end{enumerate}

% =========================================================
% ТРЕНИРОВОЧНЫЙ БЛОК
% =========================================================

\clearpage

\topic{Тренировочный блок}

\textbf{Формат: Путь А.}

Задания 1--3 --- среднего уровня,
задания 4--9 --- выше среднего,
задание 10 --- сложное.

В вычислительных заданиях записывайте основные промежуточные действия.
Калькулятор в упражнениях на технику вычислений не используйте.

\begin{enumerate}[label=\textbf{\arabic*.}]

    % -----------------------------------------------------
    % 1
    % -----------------------------------------------------

    \item
    Вычислите и запишите ответ десятичной дробью:

    \[
    \frac{7}{15}-\frac{2}{9}.
    \]

    % -----------------------------------------------------
    % 2
    % -----------------------------------------------------

    \item
    Выполните деление и запишите ответ десятичной дробью:

    \[
    \frac{14}{25}:\frac{7}{10}.
    \]

    % -----------------------------------------------------
    % 3
    % -----------------------------------------------------

    \item
    Переведите число

    \[
    3\frac25
    \]

    в следующие виды:

    \begin{enumerate}[label=\arabic*)]
        \item неправильная дробь;
        \item десятичная дробь.
    \end{enumerate}

    % -----------------------------------------------------
    % 4
    % -----------------------------------------------------

    \item
    Выполните действия столбиком:

    \begin{enumerate}[label=\arabic*)]
        \item
        \[
        2{,}34+3{,}15;
        \]

        \item
        \[
        2{,}34\cdot2{,}1.
        \]
    \end{enumerate}

    % -----------------------------------------------------
    % 5
    % -----------------------------------------------------

    \item
    Вычислите:

    \[
    4{,}2:0{,}07.
    \]

    % -----------------------------------------------------
    % 6
    % -----------------------------------------------------

    \item
    Какое из чисел

    \[
    0{,}14;
    \qquad
    0{,}18;
    \qquad
    0{,}24;
    \qquad
    0{,}30
    \]

    находится между числами

    \[
    \frac16
    \quad\text{и}\quad
    \frac15?
    \]

    % -----------------------------------------------------
    % 7
    % -----------------------------------------------------

    \item
    Между какими соседними целыми числами находится

    \[
    \sqrt{70}?
    \]

    % -----------------------------------------------------
    % 8
    % -----------------------------------------------------

    \item
    Известно, что

    \[
    -3<a<-2.
    \]

    Какое утверждение верно для любого такого $a$?

    \begin{enumerate}[label=\arabic*)]
        \item
        \[
        a+3>0;
        \]

        \item
        \[
        a+2>0;
        \]

        \item
        \[
        \frac1a>0;
        \]

        \item
        \[
        -a<0.
        \]
    \end{enumerate}

    % -----------------------------------------------------
    % 9
    % -----------------------------------------------------

    \item
    Пусть

    \[
    A=-2,
    \qquad
    B=-1,
    \qquad
    C=1.
    \]

    Найдите наибольшее из чисел:

    \[
    A+B,
    \qquad
    \frac1B,
    \qquad
    \frac AB,
    \qquad
    \frac BC.
    \]

    % -----------------------------------------------------
    % 10
    % -----------------------------------------------------

    \item
    Расположите числа в порядке возрастания:

    \[
    \frac{5}{12},
    \qquad
    0{,}44,
    \qquad
    \frac{7}{15},
    \qquad
    \frac{\sqrt2}{3}.
    \]

\end{enumerate}

% =========================================================
% ОТВЕТЫ
% =========================================================

\clearpage

\topic{Ответы к тренировочному блоку}

\renewcommand{\arraystretch}{1.4}

\begin{table}[h]
\centering
\caption{Краткие ответы}

\begin{tabularx}{\linewidth}{@{}>{\bfseries}p{12mm}X@{}}

\toprule
№ & Ответ
\\
\midrule

1 &
\[
\frac{11}{45}=0{,}2(4)
\]
\\

2 &
\[
0{,}8
\]
\\

3 &
1) $\dfrac{17}{5}$;
\qquad
2) $3{,}4$
\\

4 &
1) $5{,}49$;
\qquad
2) $4{,}914$
\\

5 &
\[
60
\]
\\

6 &
$0{,}18$, поскольку
\[
0{,}1(6)<0{,}18<0{,}2
\]
\\

7 &
\[
8<\sqrt{70}<9
\]
\\

8 &
1) $a+3>0$
\\

9 &
\[
\frac AB=2
\]
\\

10 &
\[
\frac{5}{12}
<
0{,}44
<
\frac{7}{15}
<
\frac{\sqrt2}{3}
\]
\\

\bottomrule

\end{tabularx}
\end{table}

\vspace{1em}

\textbf{Контрольные ориентиры для задания 10:}

\[
\frac{5}{12}
\approx
0{,}4167,
\]

\[
0{,}44=0{,}44,
\]

\[
\frac{7}{15}
\approx
0{,}4667,
\]

\[
\frac{\sqrt2}{3}
\approx
0{,}4714.
\]

\end{document}